\section{Implementierung der Software}
Die Softwareentwicklung erfolgt schrittweise und modular.  
Zunächst werden grundlegende Funktionen implementiert.  
Dazu zählen die Initialisierung der Hardware sowie die Ansteuerung der LED-Streifen.  
Die Ansteuerung der LEDs erfolgt über dedizierte Funktionen.  
Eine zentrale Funktion stellt die gezielte Aktivierung einzelner LEDs dar:  
\begin{verbatim}
void setLed(int ledNumber, uint8_t r, uint8_t g, uint8_t b)
<HIER KOMMT NOCH EIN CODEAUSSCHNITT HIN>
\end{verbatim}
Für die gleichzeitige Ansteuerung mehrerer LEDs werden Arrays verwendet.  
Diese enthalten die Indizes der LEDs, die ein bestimmtes Wort darstellen.  
Für Zeitanzeige, Temperaturinfos, Wetterinfis und Zusatzinformationen werden separate LED-Arrays definiert.  
Die Verarbeitung der Zeit erfolgt durch zyklisches Auslesen der Systemzeit.  
Die Anzeige wird in regelmäßigen Intervallen aktualisiert.  
Die Funktion zur Darstellung der Uhrzeit ist wie folgt aufgebaut:  
\begin{verbatim}
void showHourWord(int hour24)
<HIER KOMMT NOCH EIN CODEAUSSCHNITT HIN>
\end{verbatim}
Die Software wird iterativ aufgespielt und erweitert.  
Nach jeder Erweiterung erfolgt eine Funktionsprüfung.  
Dieses Vorgehen ermöglicht eine schrittweise Integration aller Systemfunktionen und verreinfacht die Fehlersuche bei Problemen.